%% ========================================================================
%% Brahmagupta's Brahmasphutasiddhanta Part 3 -- Chunk 2
%% Pages 1817--1956 of the original edition
%% Chapters: Triprasnottaradhyaya, Grahana-uttara-adhyaya
%% ========================================================================
\documentclass[11pt,a4paper]{article}

%% -------- Core Packages --------
\usepackage{fontspec}
\usepackage{polyglossia}
\usepackage{amsmath,amssymb,amsthm}
\usepackage{geometry}
\usepackage{graphicx}
\usepackage{xcolor}
\usepackage{titlesec}
\usepackage{fancyhdr}
\usepackage{hyperref}
\usepackage{indentfirst}

%% -------- Language Setup --------
\setmainlanguage{english}
\setotherlanguage{sanskrit}

%% -------- Page Geometry --------
\geometry{
  paperwidth=210mm,paperheight=297mm,
  left=25mm,right=25mm,top=30mm,bottom=30mm,
  headheight=14pt,footskip=30pt
}

%% -------- Font Configuration --------
\setmainfont{Linux Libertine O}[Renderer=Harfbuzz]
\setsansfont{Linux Biolinum O}[Renderer=Harfbuzz]

%% Devanagari (Sanskrit/Hindi) Font
\newfontfamily\devanagarifont[Script=Devanagari,Scale=1.1]{Noto Serif Devanagari}
\newfontfamily\devanagarifontsf[Script=Devanagari,Scale=1.1]{Noto Sans Devanagari}

%% -------- Colors --------
\definecolor{titlecolor}{RGB}{45,55,72}
\definecolor{sectioncolor}{RGB}{55,65,81}
\definecolor{notecolor}{RGB}{120,53,15}
\definecolor{rulecolor}{RGB}{180,160,140}

%% -------- Section Formatting --------
\titleformat{\section}
  {\normalfont\Large\bfseries\color{sectioncolor}}
  {\thesection}{1em}{}
\titleformat{\subsection}
  {\normalfont\large\bfseries\color{sectioncolor}}
  {\thesubsection}{1em}{}

%% -------- Header/Footer --------
\pagestyle{fancy}
\fancyhf{}
\fancyhead[L]{\small\textit{Brahmasphutasiddhanta Part 3}}
\fancyhead[R]{\small\thepage}
\renewcommand{\headrulewidth}{0.4pt}
\renewcommand{\footrulewidth}{0pt}

%% -------- Custom Commands --------
\newcommand{\longline}{\noindent\rule{\textwidth}{0.4pt}}
\newcommand{\shortline}{\noindent\rule{0.5\textwidth}{0.4pt}\par}
\newcommand{\dev}[1]{{\devanagarifont #1}}
\newcommand{\san}[1]{{\devanagarifont\textbf{#1}}}
\newcommand{\hin}[1]{{\devanagarifont #1}}

%% Verse environment for shlokas
\newenvironment{shloka}{\begin{quote}\devanagarifont}{\end{quote}}
\newenvironment{hindicom}{\devanagarifont}{\par}

%% Helper: Devanagari text outside math (avoids scriptfont issues)
\newcommand{\dtxt}[1]{\mbox{\devanagarifont #1}}

%% -------- Hyperref Setup --------
\hypersetup{
  colorlinks=true,
  linkcolor=sectioncolor,
  citecolor=sectioncolor,
  urlcolor=sectioncolor,
  pdfencoding=auto
}

%% ========================================================================

\begin{document}

%% ========================================================================
%% TITLE PAGE
%% ========================================================================
\thispagestyle{empty}
\begin{center}
\vspace*{2cm}

{\Huge\color{titlecolor}\devanagarifont\bfseries ब्राह्मस्फुटसिद्धान्त}\\[0.8cm]
{\LARGE\color{sectioncolor}\devanagarifont भाग ३}\\[1.5cm]

\shortline

{\Large\bfseries Brāhmasphuṭasiddhānta}\\[0.3cm]
{\large\itshape Part 3 --- Pages 1817--1956}\\[1cm]

\shortline

{\large By}\\[0.3cm]
{\LARGE\bfseries Brahmagupta}\\[0.5cm]
{\large (b.~598 CE)}\\[1.5cm]

\shortline

{\large Edited with Sanskrit Commentary (Vāsanā) and Hindi Translation}\\[0.3cm]
{\large by}\\[0.3cm]
{\large\bfseries Sudhākara Dvivedī}\\[1cm]

\shortline

{\small This digital edition covers original pages 1817--1956, comprising:}\\[0.2cm]
{\small\devanagarifont त्रिप्रश्नोत्तराध्यायः (Tri-prashnottara-adhyāya)}\\[0.1cm]
{\small\devanagarifont ग्रहणोत्तराध्यायः (Grahana-uttara-adhyāya)}\\[0.1cm]
{\small Editorial Introduction to the chapters}\\[2cm]

\vfill
{\small Typeset in XeLaTeX with Noto Serif Devanagari}\\
{\small \today}
\end{center}
\newpage

%% ========================================================================
%% TABLE OF CONTENTS
%% ========================================================================
\tableofcontents
\newpage

%% ========================================================================
%% CHAPTER XV: TRIPRASNOTTARADHYAYA
%% ========================================================================
\section*{\dev{त्रिप्रश्नोत्तराध्यायः}}
\addcontentsline{toc}{section}{त्रिप्रश्नोत्तराध्यायः (Chapter XV: Tri-prashnottara-adhyāya)}
\markboth{Tri-prashnottara-adhyāya}{Tri-prashnottara-adhyāya}

\noindent\rule{\textwidth}{0.6pt}
\begin{center}
\textit{Original pages 1033--1073}
\end{center}
\noindent\rule{\textwidth}{0.6pt}
\vspace{0.5cm}

\begin{hindicom}
\dev{अब श्रम्बिज्ञुत नक्षत्र के भोगानयन और ग्रहमुक्त नक्षत्रानयन को कहते हैं।}
\end{hindicom}

\vspace{0.3cm}
\begin{shloka}
\dev{पञ्चदशो योगभोगकलानामैक्यं = चं × १५}

\dev{सर्वेषां योगः = १ चं + ३ चं + १५ चं = २७ चं = सर्वनक्षत्रभोगः}
\end{shloka}

\begin{hindicom}
\dev{चक्कलास्यः शुद्धाः सर्वनक्षत्रभोगसंख्यास्तदा श्रम्बिज्ञुद् भोगकलास्तदिह्नगतिः =}

\dev{चक्र—२७ चं ततः कल्पोड्भवो भगणः := (चक्र—२७ चं) कुदिन =}

\dev{चक्कला}

\dev{कुदिन—२७ कल्पचंभगणा एतेनाचार्येण तुमूपपन्नम् । सिद्धान्त शेखरे}

\dev{``चक्कारिया वा शाश्वतसतिवधनाहतानि शोधयिन् सुदिनलयद्वरोषचक्र'' । स्यादे-}
\dev{कवासरम्बा कलिकागतिर्या सा वैवस्वेतावमध्यगाभिध्या मृतिः ॥ इद्धप्रहस्य}
\dev{कलिकालिकरादिसंशया नक्षत्रभोगकलिका प्रहसुत्कालानि । शेषात् भोग्यकलिका}
\dev{परितात् गम्य तास्यो भवन्ति गतगम्यदिनानि मुक्त्या ॥ इति श्रीपतेः पञ्चद्वय}
\dev{``सर्वेषां भोगैनितचक्कलिप्ता वैवस्वतः स्यादमिज्ञुद्भयोग'' इत्यादि भार्करोक्तं}
\dev{च सर्वाचार्येणस्य सर्वेषं समानार्थकार्मित ॥ ५०--५१ ॥}
\end{hindicom}

\vspace{0.5cm}
\noindent\textbf{\dev{उपपत्तिः}}

\begin{shloka}
\dev{षट् श्रद्धं भोगकला नक्षत्रों का ऐक्य =} $\frac{3\ \dtxt{चं}}{2} \times \dtxt{६} = ३$ \dev{चं × ३ = ६ चं}

\dev{षट् शद्धभोग कला नक्षत्रों का ऐक्य =} $\frac{\dtxt{चं}}{2} \times \dtxt{६} =$ \dev{चं × ३}

\dev{पनद्वह् एक भोग (समान भोग) कलानक्षत्रों का ऐक्य = चं × १५}

\dev{सबों का योग = सर्वनक्षत्र भोग कला = ६ चं + ३ चं + ? चं = २७ चं}
\end{shloka}

\vspace{0.5cm}
\begin{hindicom}
\dev{हिं. भा.---चन्द्रमध्यगति कला को ईं, ईं, ? इन से गुणा करने से षड्वर्षं, शर्बं,}
\dev{सम नक्षत्रों की भोग कला होती है । सब नक्षत्रों की जो भोग कला है उनको भगण कला}
\dev{में से घटाने से जो शेष रहता है वह श्रम्बिज्ञुत भोग है । वा कल्प चन्द्र भगण सतासी से}
\dev{गुणा कर कल्पकुदिन में से घटाने से शेष श्रम्बिज्ञुत का कल्प भगण होता है, इन भगणों से}
\dev{जो दिन भोग (कल्प कुदिन में यदि श्रम्बिज्ञुत का कल्प भगण पाते हैं तो एक दिन में क्या इस}
\dev{अनुपात से लब्ध कालात्मक दिनगति) होता है वह श्रम्बिज्ञुत भोग होता है । ग्रहकालसमूह}
\dev{में से नक्षत्र भोग कला को घटाने से नक्षत्र होते हैं, अर्थाद् ग्रहमुक्त कला में जितने नक्षत्रों की}
\dev{भोग कला शुद्ध हो उतने गतनक्षत्र होते हैं, शेषकला वर्तमान नक्षत्र की गत कला है उसको}
\dev{नक्षत्र भोग कला में से घटाने से गम्य कला होती है । तब ग्रहगति कला में एक दिन पाते हैं}
\dev{तो गत-गम्य कला में से क्या इस अनुपात से गतदिन और गम्यदिन होते हैं ॥ ५०--५२ ॥}
\end{hindicom}

\vspace{0.5cm}
\noindent\textbf{\dev{उपपत्तिः}} --- \dev{यहाँ आचार्य ने पहले उक्त कालज्ञा को इष्टद्वित कल्पित किया है । तब श्रदा,}
\dev{समशाईं कुदूदित् इन तीनों मुञ्जाओं से उत्पन्न एक श्रम्ब क्षेत्र है, तथा श्रदज्ञा, लब्धज्ञा-}
\dev{तिज्ञा इन तीनों मुञ्जाओं से उत्पन्न द्वितीय श्रम्ब क्षेत्र है, ये दोनों श्रम्ब क्षेत्र सजातीय हैं}
\dev{इसलिए अनुपात करते हैं}

\begin{equation*}
\frac{\dtxt{लब्धज्ञा. इष्ट}}{\dtxt{तिज्ञा}} = \dtxt{समशाईंकुद्}
\end{equation*}

\dev{तथा कालज्ञा =}

\vspace{0.5cm}

\subsection*{\dev{अत्रोपपत्तिः}}

\begin{shloka}
\dev{छायाकर्णे गोले पलभा शङ्कुतुल्योस्तुल्यतत्त्वं भवति कथमिति प्रदर्शयते यदि}
\dev{द्वादशाङ्कुलराश्यों पलभा भुजां लभ्यते तदेप्तशाङ्कुं किमिति जातं शङ्कुतुल्यम् =}
\dev{$\frac{\dtxt{पमा. शाङ्कु}}{१२}$, परन्तु शाङ्कु = $\frac{१२ \times \dtxt{त्र}}{\dtxt{छाक}}$ शत उत्थापनेन शङ्कुतुल्यम् = $\frac{\dtxt{पमा. १२. त्र}}{१२.\ \dtxt{छाक}}$}
\end{shloka}

\vspace{0.3cm}
\dev{यहाँ द्वाद (दिनार्क्) से भिन्न समय में पलभा साधन के लिए कहते हैं ।}

\dev{हिं. भा.---दिनार्क् से भिन्न समय में रवि की छाया को छायाकर्णों से गुणा कर}
\dev{तिज्ञा से भाग देने से छायाकर्ण गोलिय छाया होती है, शङ्कुमूल और पूर्वापर रेखा का}
\dev{भेद भुज है । उत्तर गोल में कर्णदुताज्ञा में उत्तर भुज को घटाने से और दक्षिणा भुज}
\dev{को जोड़ने से पलभा होती है । दक्षिणा गोल में भुज में कर्णदुतिय छाया को घटाने ही से}
\dev{पलभा होती इति ॥ ४६--५० ॥}

\newpage

%% ========================================================================
\section*{\dev{ग्रहणोत्तराध्यायः}}
\addcontentsline{toc}{section}{ग्रहणोत्तराध्यायः (Chapter XVI: Grahana-uttara-adhyāya)}
\markboth{Grahana-uttara-adhyāya}{Grahana-uttara-adhyāya}

\noindent\rule{\textwidth}{0.6pt}
\begin{center}
\textit{Original pages 1089--1139}
\end{center}
\noindent\rule{\textwidth}{0.6pt}
\vspace{0.5cm}

\begin{hindicom}
\dev{अब श्रम्ब्य प्रश्नों को कहते हैं ।}

\dev{हिं. भा.---जो क्रान्ति के ज्ञाता समशकु को जानते हैं । जो लम्बांश वेत्सा समशुत्त}
\dev{को जानते हैं, कोणशकु को जानते हैं, कोणशकुछाया को जानते हैं । कोणदुताज प्रवेश में}
\dev{घटी को जानते हैं वे ज्योतिष सिद्धान्त के पण्डित हैं, यहाँ पाँच प्रश्न हैं ॥ ७ ॥}
\end{hindicom}

\vspace{0.5cm}
\noindent\textbf{\dev{इदानीमन्यान् प्रश्नानाह}}

\begin{shloka}
\dev{शङ्कुतलप्राच्यपरान्तरद्वयं बोध्यं यो विजानाति ।}
\dev{विशुवच्छायामेकं हस्त्वाऽऽसन्नदित्यं च गण्यकः सः ॥ ७ ॥}
\end{shloka}

\vspace{0.3cm}
\begin{hindicom}
\dev{सू. भा.---शङ्कु तलप्राच्यपरान्तरं भुजः । यो भुजद्वयं बोध्य विशुवच्छाया}
\dev{पलभां विजानाति । एकमेव भुजं हस्त् वा विशुवच्छायामादित्यमकं च विजानाति}
\dev{स गण्यकः । एवमत्र प्रश्नत्रयम ॥ ८ ॥}
\end{hindicom}

\vspace{0.3cm}
\begin{hindicom}
\dev{वि. भा.---शङ्कुतलपूर्वापररेखयोरन्तरं भुजोर्धित, यो भुजद्वयं ज्ञात्वा पलभां}
\dev{जानाति एकं भुजं ज्ञात्वा पलभां रविं च जानाति स गण्यकोऽस्तीति । अत्र प्रश्नत्रय-}
\dev{भमस्ति ॥ ७ ॥}
\end{hindicom}

\vspace{0.5cm}
\begin{hindicom}
\dev{अब श्रम्ब्य प्रश्नों को कहते हैं ।}

\dev{हिं. भा.---शङ्कुतल और पूर्वापर रेखा का भेद भुज है । जो व्यक्ति दो भुजों को}
\dev{ज्ञान कर पलभा को जानते हैं । एवं एक भुज को ज्ञान कर पलभा को जानते हैं तथा रवि}
\dev{को जानते हैं वे ज्योतिः शास्त्र के पण्डित हैं इति । यहाँ तीन प्रश्न हैं ॥ ७ ॥}
\end{hindicom}

\vspace{0.5cm}
\noindent\textbf{\dev{इदानीमन्यान् प्रश्नानाह}}

\begin{shloka}
\dev{पातालशङ्कुमुध्येस्ते वा हस्तज्ञां रविंविजानाति ।}
\dev{इष्टपातालगणकोः पृथक् तले वा स तन्त्रज्ञः ॥ ८ ॥}
\end{shloka}

\vspace{0.3cm}
\begin{hindicom}
\dev{सू. भा.---यो रवेः पातालशङ्कु कुमथः शङ्कुं विजानाति । उदयेस्ते वा यो}
\dev{रविंद्‌र्ग् ज्ञात्रो विजानाति । इष्टशङ्कुकुदिवोइष्टशङ्कुः पातालगणःः पातालशङ्कुः}
\dev{रविरथः शङ्कुः । तयोः पृथक् पृथक् तले शङ्कुतले च वा यो विजानाति स एव}
\dev{तन्त्रज्ञ इति । एवमत्र प्रश्नचतुष्टयम ॥ ९ ॥}
\end{hindicom}

\vspace{0.5cm}
\begin{hindicom}
\dev{अब श्रम्ब्य प्रश्नों को कहते हैं ।}

\dev{हिं. भा.---जो व्यक्ति ग्रहण में राहुमार्ग (भूमामार्ग) को जानते हैं । उस मार्ग}
\dev{को जानते हैं, कोणशकुछाया को जानते हैं । यहाँ पाँच प्रश्न हैं ॥ ७ ॥}
\end{hindicom}

\vspace{0.5cm}

\begin{shloka}
\dev{कोणदुतस्थरविगतं भुवप्रोत्कुटं भुवाद्रिविपर्यं भुव्यचापांशं एकं भुजः ।}
\dev{भुवाल्पसतिकावधि लम्बांश द्वितीयो भुजः । कोणदुतं व स्वैनिकाद्रिविपर्यं ननां-}
\dev{शास्तुतीयो भुजः । त्रिभुजेऽस्मिन् रविगतं भुवप्रोत्कुटतयाप्येतर्कुर्नाम्यामृत्यन्न-}
\dev{कोणा नतकालः । याम्योत रवित्कोणदुतनाम्यामृत्यन्नकोणः = ८५ नतोनुपातः}
\dev{किम्यते यदि नतकालज्ञया हस्तज्ञा लभ्यते तदाश्रवेदांशज्ञया (ज्ञा ८५) किमिनि}
\dev{समागच्छति भुज्या = $\frac{\dtxt{हस्तज्ञा. ज्ञा ८५}}{\dtxt{नतकालज्ञा}} = \frac{\dtxt{हस्तज्ञा. त्रि. ज्ञा ८५}}{\dtxt{नतकालज्ञा. त्रि}} = \frac{\dtxt{हस्तज्ञा. त्रि. २८५}}{२४९५\ \dtxt{नतकालज्ञा}}$}
\end{shloka}

\newpage

\subsection*{\dev{उपपत्तिः}}

\begin{hindicom}
\dev{यहाँ देशान्तर विदित नहीं है इसलिए देशान्तर संस्कार से रहित स्फुट रवि और}
\dev{स्फुट चन्द्र से चन्द्र ग्रहण विधि से सर्वग्रहण में संमीलन काल और उन्मीलन काल साधन}
\dev{करना चाहिए । उस दिन में द्वितीय से भी संमीलन काल समज्ञा चाहिए वह यदि गणितानगत्}
\dev{संमीलन काल से अधिक हो तो द्वितीय रेखा से पूर्व भाग में अथवा पश्चिम भाग में स्थित है यह}
\dev{समज्ञा चाहिए । क्योंकि रेखा से पूर्व स्वदेश में पहले स्वदेश में पश्चात् रेखा देश में मध्याह्न}
\dev{काल होता है । अतः रेखा देशीय इष्ट संमीलन काल से स्वदेशीय संमीलन काल अधिक होता}
\dev{है पश्चिम में इसके विपरीत होता है । इस तरह उन्मीलन काल से इष्ट प्रास काल से स्पर्श}
\dev{काल से या मोक्ष काल से परीक्षा होती है, स्पर्श काल परीक्षा और मोक्ष काल परीक्षा इष्ट}
\end{hindicom}

\vspace{0.5cm}
\begin{shloka}
\dev{उत्तरगोले याम्ये विषुवच्छायायुक्तस्तरं हीनम् ।}
\dev{एवं विषुवच्छाया युक्तविहीनास्तरेषाभा ॥ ५० ॥}
\end{shloka}

\begin{hindicom}
\dev{सू. भा.---प्रथमायाः पूर्वार्धं त्रिप्रश्नाध्यायस्य चतुर्थीयुपयुपर्धसंं न्यस्यात्-}
\dev{मेव । अन्यत् सर्वं च त्रिप्रश्ना ५८--६० आयामिः स्फुटम् ॥ ४९--५० ॥}
\end{hindicom}

\begin{hindicom}
\dev{वि. भा.---दिनार्धदिनकाले शक्रिया छायाकर्णगुणा तिज्ञामक्ता तदा}
\dev{छायाकर्णगोलिया रवैर्या भवेत् । शङ्कुमूलपूर्वापररेखयोरन्तरं भुजः । तेन भुजे-}
\dev{नोन्यता कर्णदुताप्राच्यादितुरेया भुजेनोनदिक्षेत्रं भुजेन युता तदोतरगोले}
\dev{विषुवच्छाया (पलभा) भवेत् । याम्ये (दक्षिणा गोले) तया कर्णदुतियाप्राच्यादन्तरं}
\dev{(भुजमान) हीन तदा पलभा भवेत् । एवमन्तरेषा (भुजेन) युक्तविहीनास्तस्यथा-}
\dev{दन्तरभुजेन युता दक्षिणाम्यभुजेन हीना तदा पलभा भवेदिति ॥ ४९--५० ॥}
\end{hindicom}

\vspace{0.5cm}
\noindent\textbf{\dev{अत्रोपपत्तिः}}

\begin{hindicom}
\dev{छायाकर्ण गोले पलभा शङ्कु तुल्ययोस्तुल्यतत्त्वं भवति कथमिति प्रदर्शयते यदि}
\dev{द्वादशाङ्कुलराश्यों पलभा भुजां लभ्यते तदेप्तशाङ्कुं किमिति जातं शङ्कुतुल्यम् =}

\begin{equation*}
\frac{\dtxt{पमा. शाङ्कु}}{१२}, \quad \dtxt{परन्तु शाङ्कु} = \frac{१२ \times \dtxt{त्र}}{\dtxt{छाक}} \quad \dtxt{शत उत्थापनेन}
\end{equation*}

\begin{equation*}
\dtxt{शङ्कुतुल्यम्} = \frac{\dtxt{पमा. १२. त्र}}{१२.\ \dtxt{छाक}} = \dtxt{पमा}
\end{equation*}
\end{hindicom}

\vspace{0.3cm}
\begin{hindicom}
\dev{$\therefore$ सिद्धं कर्णगोले शङ्कुतुल्यम् = पलभा । कर्णदुताज्ञा व्यस्तगोला भवति ।}
\dev{पलभा च सव्योतरा । तथाः संस्कारतःछायाप्रविपरसूनमध्ये भुजः कथतेऽस्तत्-}
\dev{क्षेपरियेन कर्णदुताज्ञा भुजयाः संस्करेणा पलभा भवतीति । सिद्धान्तशेखरे ``अथ-}
\dev{या तु नतपूर्वपश्चिमाशान्तरेषा शकिबुतज्ञामका । दक्षिणोत्तरभुजा युतोनिता}
\dev{सौम्यवर्तिन् रवो पलभा ॥ दक्षिणेन पुनरिननायां तथा हीनमेव हि तदन्तरं}
\dev{सदा । एवमन्तरयुतोनिता भवेद्-छाया नियतमुगुलाश्चका ॥'' श्रीपत्युक्तसिद्धमा-}
\dev{चार्येणानुरूपमेवेति ॥ ४९--५० ॥}
\end{hindicom}

\vspace{0.5cm}
\begin{shloka}
\dev{वलनस्यमुचे तैम्यः पृथग् ग्रहक् खण्डकेन । वृत्तिः कृतिः स्पर्शविमर्शनिमध्यग्रता}
\dev{क्रमेणोर्बिमहावगमयः ॥}
\end{shloka}

\begin{hindicom}
\dev{``भास्कराचार्येण श्रीपतिप्रकार एव विश्वदर्शनेण प्रति}
\dev{पादितः ॥ १४--१५॥}
\end{hindicom}

\vspace{0.5cm}
\begin{hindicom}
\dev{विषुवदर्पमण्डलदिका इत्यादि दो प्रश्नों के उत्तर ग्रहणाधिकार में बता दिए}
\dev{गये हैं । उसके उत्तरार्ध को श्रम्बे कहते हैं ।}

\dev{अब `सम्पर्कं मण्डलै' इत्यादि प्रश्न के उत्तर कहते हैं ।}

\dev{हिं. भा.---प्राह्लादूता में स्वल्पान्तर से सरलाकार तात्कालिक क्रान्तिज्ञुनीय चाप}
\dev{बलन सूत्र का । बलनज्ञा से दिव्यज्ञान समज्ञा चाहिए । मानचक्रबाहृत में बलन सूत्र के}
\dev{उपर रवि के स्पर्श कालिक सार और मोक्ष कालिक सार को जिस दिशा के सार है उसी दिशा}
\dev{में देना । मध्यबलन सूत्र में प्राह्ल बिम्ब केन्द्र से मध्यार देना चाहिए । चन्द्र शराघ्रों से प्राह्ल}
\dev{प्रमाण व्यास से तुरत लिखकर स्पर्श मोक्ष श्रीर प्रास समज्ञा चाहिए, एवं पृथिवी पर परिलेख}
\dev{लिखने से स्पर्श मोक्ष श्रीर प्रास होता है इति ॥}
\end{hindicom}

\vspace{0.5cm}
\noindent\textbf{\dev{उपपत्तिः}}

\begin{hindicom}
\dev{मानचक्रबाहृत् वृत्त में जब प्राह्लाकृता का केन्द्र होता है तब प्राह्ल बिम्बग्रान्त और प्राह्ल}
\dev{बिमग्रान्त संलग्न रहता है । इस लिये विदित मानचक्रबाहृत् वृत्त में दिशा को श्रीष्ट करना}
\dev{चाहिए । सममण्डल और मानचक्रबाहृत् वृत्त का समाप्त बिन्दु मानचक्रबाहृत् वृत्त में प्राची चिह्न है ।}
\dev{वहाँ से बलनज्ञा दान देने से केन्द्र से बलनज्ञाप्रगत रेखा क्रान्तितृत् प्राची है बलनसूत्र से}
\dev{ज्यावत् खादान देना चाहिए । क्योंकि क्रान्तिबृत् प्राची से सार दक्षिण और उत्तर रहता है । इस}
\dev{तरह स्पर्श और मोक्ष में होता है । मध्यबलन तात्कालिकक्रान्तिबृत् प्राची से दक्षिणोत्तर}
\dev{दिशा में होता है इसलिए केन्द्र से बलन सूत्र में देना चाहिए । शराघ्र में प्राह्ल वृत्त का केन्द्र}
\dev{होता है इसलिए वहाँ से तात्कालिकप्राह्लाभि प्रमाण से रचित बृत्तों से स्पर्श मोक्षमध्य होने है}
\dev{सिद्धान्त शेखर में `मानचक्रबाहृतलये तत्परियामिनिवि' इत्यादि संस्कृतोपपत्ति में लिखित}
\dev{ल्लोक से श्रीपति ने आचार्यों के अनुकूल ही कहा है । सिद्धान्त शिरोमणि में `प्राह्लाखसूत्रेय}
\dev{विधाय वृत्तं मानचक्रबाहृतं च साधितांशम्' इत्यादि से भास्कराचार्य ने श्रीपति प्रकार ही को}
\dev{विचार रूप से प्रतिपादित किया है इति ॥ १४--१५ ॥}
\end{hindicom}

\vspace{0.5cm}
\begin{shloka}
\dev{इदानी यः परिलिखतीष्टद्यासमिल्यादि प्रश्नस्योत्तरमाह ।}
\dev{पश्चात् प्रग्रहयो प्रागमोषे रविविब्बलो बाहुः ।}
\dev{स्वबलनसिद्धज्ञादिशि विपरीतः शीतकरमध्यात् ॥ १६ ॥}
\dev{भानुमतो बाहुगुण्याख्या दिशं कोटिरत्यथा शविनः ।}
\dev{रविशास्तिवमध्यात् करणास्तिन्यकं, करणांगुकोटियुतेः ॥ १७ ॥}
\end{shloka}

\newpage

%% ========================================================================
\section*{\dev{सूर्यसिद्धान्तयन्त्रविवरण}}
\addcontentsline{toc}{section}{सूर्यसिद्धान्तयन्त्रविवरण (Surya-siddhanta Yantra-vivarana)}
\markboth{Surya-siddhanta Yantra-vivarana}{Surya-siddhanta Yantra-vivarana}

\noindent\rule{\textwidth}{0.6pt}
\begin{center}
\textit{Original pages 1121--1139}
\end{center}
\noindent\rule{\textwidth}{0.6pt}
\vspace{0.5cm}

\begin{hindicom}
\dev{सूर्य सिद्धान्त में शब्दों लिखित यन्त्र विवरण है---}
\end{hindicom}

\begin{shloka}
\dev{तुङ्गार्कसमायुक्तं गोलयन्त्रं प्रसाधयेत् ।}
\dev{गोपयेेतत्प्रकाशोक्तं सर्वगंयं भवेद्धि ॥}

\dev{कालसंसाधनार्थं तथा यन्त्रारिणा साधयेत् ।}
\dev{एकाकी योजयेद्‌बीजं यन्त्रं विस्मय कारिणि ॥}

\dev{शङ्कुकुयष्टि धनुश्चक्रेषछायायन्त्रैर्नैकथा ।}
\dev{गुप्तदेशाङ्कज्ञेयं कालज्ञानमतन्द्रितः ॥}

\dev{तोययन्त्रकपालाङ्कघर्म्म्युदनत्रवानरैः ।}
\dev{ससूत्ररेखुगमैश्च समयकालं प्रसाधयेत् ॥}

\dev{परदारामनुष्याणां गुल्वतैलजलानि च ।}
\dev{बीजानि पांसवस्तेषु प्रयोगास्तेऽपि दुर्लभाः ॥}

\dev{ता प्रपञ्चमधिश्छद्द् न्यस्तं कुण्डे स्मलाभमिस् ।}
\dev{द्विट्ठमेजलयहोरात्रं स्फुटं यन्त्रं कपालकम् ॥}

\dev{नयन्त्रं तथा साधु दिवा च विमले रवौ ।}
\dev{छाया संसाधनैः प्रोक्तं कालसाधनमुत्तमम् ॥}
\end{shloka}

\vspace{0.5cm}
\noindent\textbf{\dev{मानाध्याय}}

\begin{shloka}
\dev{मानानि सौरचान्द्रार्क्षसावनानि ग्रहानयनमैभिः ।}
\dev{मानैः पृथक् चतुर्भिः संख्यबहारोड्ज्ञ लोकस्य ॥}
\end{shloka}

\begin{hindicom}
\dev{इससे सौरमान, चान्द्रमान, नाक्षत्रमान और सावनमान, ये चार प्रकार के मान}
\dev{कहे गये हैं । इन्हीं चारों मानों से लोगों के सब व्यवहार होते हैं । किस किस मान से कौन}
\dev{कौन पदार्थ ग्रहण किये जाते हैं यह सब प्रति पादित है । बाह्म, दिव्य, पिट्र, प्राजापत्य,}
\dev{बाहृस्पत्य, सौर, सावन, चान्द्र, नाक्ष ये नौ मान हैं । इन मानों में से मनुष्यलोक}
\dev{में केवल सौर, चान्द्र, सावन और नाक्ष इन चार मानों की ही प्रधानता है । क्योंकि}
\dev{इन्हीं मानों से मनुष्यों के सब व्यवहार सम्पन्न होते हैं । सूर्यसिद्धान्त, सिद्धान्तशेखर, सिद्धान्त-}
\dev{शिरोमणि आदि सब ग्रन्थों में मानों के विषय में समान रूप से कहा गया है । ब्रह्मस्फुट}
\dev{सिद्धान्त के इस अध्याय में भूमादर्शन के भी साधन हैं ।}
\end{hindicom}

\newpage

%% ========================================================================
\section*{\dev{स्मृत्योर्युगलयोः समन्वयः}}
\addcontentsline{toc}{section}{स्मृत्योर्युगलयः समन्वय (Reconciliation of Smrti Texts)}
\markboth{Smriti Reconciliation}{Smriti Reconciliation}

\noindent\rule{\textwidth}{0.6pt}
\begin{center}
\textit{Original pages 1133--1139}
\end{center}
\noindent\rule{\textwidth}{0.6pt}
\vspace{0.5cm}

\begin{hindicom}
\dev{स्मृतियों में और पुराणों में राहुग्रस्त ग्रहण का प्रतिपादन विद्यमान है । अतः दोनों मतों}
\dev{का समन्वय करते हुए आचार्य ने कहा है---}
\end{hindicom}

\begin{shloka}
\dev{राहुस्तच्छदयति प्रविशति यच्छुकलपच्छदयन्ते ।}
\dev{छाया तमसीद्वैरेघ्रदानां कमलयोनिः ॥}

\dev{चन्द्रोऽस्मुयोजः स्यो यदिनिमयभास्करस्य मासात् ।}
\dev{छादयति शर्मिततापो राहुश्छदयति तत् सवितुः ॥}
\end{shloka}

\begin{hindicom}
\dev{सिद्धान्तशिरोमणि के गोलाध्याय में भी अप्रोक्तित भास्करोक्ति---}
\end{hindicom}

\begin{shloka}
\dev{द्विदेशकालवरयाणि मेदानां छादको राहुरिति स वर्णित ।}
\dev{यन्मानिनः केवलगोलविधास्तत्सहिता वेद्यरणाबाह्याम् ॥}

\dev{राहुः क्षुभामण्डलगः शशाङ्कं शशाङ्कश्छादयतीव बिम्बम् ।}
\dev{तमोमयः शम्युवरप्रदानात् सर्वगमानामविरुद्धमेतत् ॥}
\end{shloka}

\begin{hindicom}
\dev{से समन्वय किया गया है । सिद्धान्तशेखर में राहुग्रस्त ग्रहण के खण्डनार्थ `राहुनिरा-}
\dev{करणाध्याय' नाम का एक अध्याय रखा गया है । इसमें श्रीपति ने भी निम्नलिखित}
\dev{श्लोक ने समन्वय किया है---}
\end{hindicom}

\begin{shloka}
\dev{विष्णुलूनविशिरसः किल पञ्चदेत्तवान् वरमिमं परमेष्ठी ।}
\dev{हेमदानविधिना तवतुपतिस्तमरतिमहस्तदुपरागे ॥}

\dev{भूमेच्छाया प्रविष्टः स्थगयति शशिनं शुक्लपक्षावसाने ।}
\dev{राहुर्भुप्रसादात् समविगतवस्तुतामो व्यसुतुल्यः ॥}

\dev{उर्व्यां स्थं भानुबिम्बं सलिलमयतनोर्यथैवोक्तं बिम्बम् ।}
\dev{संसूयैवं च मासव्युपरतिसमये स्वस्य साहित्यहेतोः ॥}
\end{shloka}

\begin{hindicom}
\dev{गोलबन्धाधिकार में महत् द्वृतों (पूर्वापरदृत, याम्योत्तरदृत, स्थितिजदृत आदि) की रचना}
\dev{तथा लघुदृतों (मेषादिक द्वादश राशियों के बहिरोराजदृत) की रचना करके परमलम्बन-नति}
\dev{का स्वरूप प्रतिपादन कर आचार्य ने द्वृतों का ज्ञानयन किया}
\dev{है । ब्रह्मस्फुटसिद्धान्त में आचार्य ने जैसे गोलबन्ध कहा है}
\dev{वैसे ही सिद्धान्त शेखर में श्रीपति और सिद्धान्त शिरोमणि के}
\dev{गोलाध्याय में भास्कराचार्य ने कहा है । ग्रहगोल और नक्षत्रगोल में पाँच स्थिरदृत}
\dev{(पूर्वापरदृत, स्थितिजदृत, याम्योत्तरदृत, उन्मण्डल, विषुवदृत) कहे हैं । ये सब कला-}
\dev{मण्डल के बराबर हैं । तथा ग्रहों के चलदृत मन्दनीचोचदृत = ७, मन्दादि ग्रहों के शीघ्र-}
\dev{नीचोचदृत = ५ । मन्दप्रतिबृत = ७, शीघ्रप्रतिबृत = ७ । सात ग्रहों के हमण्डल हृद्क्षेप}
\end{hindicom}

\vspace{1cm}

%% ========================================================================
%% END MATTER
%% ========================================================================
\newpage
\section*{Colophon}

\begin{center}
\shortline
\vspace{0.5cm}

{\large\textit{End of Chunk 2 (Pages 1817--1956)}}\\[0.3cm]
{\small Comprising the Tri-prashnottara-adhyāya and Grahana-uttara-adhyāya}\\[0.3cm]
{\small with editorial notes on the Sūrya-siddhānta and Smṛti reconciliation}\\[0.5cm]

\shortline

\vspace{1cm}

{\small\textit{``Brahmasphutasiddhanta'' is one of the most important works}}\\
{\small\textit{of ancient Indian mathematics and astronomy, composed by}}\\
{\small\textit{Br